\documentclass[a4paper,12pt,oneside]{article}

\usepackage[T2A]{fontenc}
\usepackage[utf8]{inputenc}
\usepackage[russian]{babel}

\usepackage{microtype}

\usepackage{amsmath,amssymb,amsthm,mathtools}

\usepackage{geometry}
\geometry{top=25mm,left=25mm,right=20mm,bottom=25mm}

\usepackage{enumitem}
\setlist[enumerate]{leftmargin=*, itemsep=0.55\baselineskip, topsep=0.35\baselineskip}
\setlist[itemize]{leftmargin=*, itemsep=0.25\baselineskip, topsep=0.25\baselineskip}

\setlength{\parindent}{1.25em}
\setlength{\parskip}{0pt}

\newcommand{\Z}{\mathbb{Z}}

\newcommand{\Task}[1]{\section*{\normalfont\bfseries Задание #1}}

\begin{document}

\begin{center}
    {\Large\bfseries Домашняя работа №2. Бинарные отношения.}\\[2mm]
\end{center}
\vspace{2mm}
\hrule
\vspace{4mm}

\Task{1}
Пусть $A=\{1,2,3,4,5\}$, $B=\{6,7,8,9\}$, $C=\{10,11,12,13\}$.\\
$R\subseteq A\times B$, $S\subseteq B\times C$, где
\[
R=\{(1,7),(4,6),(5,6),(2,8)\},\qquad
S=\{(6,10),(6,11),(7,10),(8,13)\}.
\]
\begin{enumerate}[label=\arabic*)]
\item Представьте $R$ в виде ориентированного графа (считаем, что $R\subseteq (A\times B)^2$) и в виде матрицы.
\item Определить отношения:
\begin{enumerate}[label=\alph*)]
\item $R^{-1}$;
\item $S^{-1}$;
\item $S^{-1};R^{-1}$;
\item $(R;S)^{-1}$.
\end{enumerate}
\end{enumerate}

\Task{2}
Определим $R^{-1}$ следующим образом: если $xRy$, то $yR^{-1}x$.\\
Выполнено ли соотношение $R;R^{-1}=id_A$, где $id_A$ – тождественное отношение?

\Task{3}
Докажите, что
\[
(R;S)^{-1}=S^{-1};R^{-1}.
\]

\Task{4}
Пусть $R$ и $S$ – антисимметричные отношения на $A$.\\
Будет ли антисимметричным их
\begin{enumerate}[label=\alph*)]
\item объединение?
\item пересечение?
\end{enumerate}

\subsection*{ВНИМАНИЕ}
В заданиях на свойства БО обоснование ответа обязательно. Например, для доказательства симметричности необходимо рассмотреть все пары (или воспользоваться свойствами математических операций), а для доказательства транзитивности необходимо рассмотреть все тройки (или воспользоваться свойствами математических операций), для опровержения – привести контрпример.

\Task{5}
Пусть $A=\{a,b,c,d,e\}$, а $S$, $T$, $U$ и $V$ – отношения на $A$, где
\[
S=\{(a,a),(a,b),(b,c),(b,d),(c,e),(e,d),(c,a)\},
\]
\[
T=\{(a,b),(b,a),(b,c),(b,d),(e,e),(d,e),(c,b)\},
\]
\[
U=\{(a,b),(a,a),(b,c),(b,b),(e,e),(b,a),(c,b),(c,c),(d,d),(a,c),(c,a)\},
\]
\[
V=\{(a,b),(b,c),(b,b),(e,e),(b,a),(c,b),(d,d),(a,c),(c,a)\}.
\]
Определите свойства отношений.

\Task{6}
Определите, какие из приведенных ниже отношений на $\Z^2$ являются рефлексивными, симметричными, а какие транзитивными:
\begin{enumerate}[label=\arabic*)]
\item $x+y$ – нечётное число;
\item $x+y$ – чётное число;
\item $xy$ – нечётное число;
\item $x+xy$ – чётное число.
\end{enumerate}

\Task{7}
Перечислите упорядоченные пары, принадлежащие отношениям, заданным на множестве
\[
\{x:\ x\in \Z\ \text{и}\ 1\le x\le 12\}:
\]
\begin{enumerate}[label=\arabic*)]
\item $R=\{(x,y):\ xy=9\}$;
\item $S=\{(x,y):\ 2x=3y\}$;
\item определите свойства $R$ и $S$;
\item замыкание $R$ по транзитивности;
\item замыкание $S$ по транзитивности.
\end{enumerate}

\Task{8}
Найдите замыкание по рефлексивности, по симметричности и по транзитивности отношения
\[
\{(a,a),(b,b),(c,c),(a,c),(a,d),(b,d),(c,a),(d,a)\},
\]
заданного на множестве $\{a,b,c,d\}$.\\
Имеет ли смысл строить замыкание по антисимметричности?

\Task{9}
Отношение $R$ на множестве $\Z^2$ определяется так:
\[
xRy \iff x^2-y^2\ \text{делится на}\ 3.
\]
Докажите, что $R$ – отношение эквивалентности, и опишите классы эквивалентности.

\Task{10}
Prove or disprove the following statements about the functions $f$ and $g$:
\begin{enumerate}[label=\alph*)]
\item If $f$ and $f\circ g$ are injections, then $g$ is also an injection.
\item If $f$ and $f\circ g$ are surjections, then $g$ is also a surjection.
\end{enumerate}

\Task{11}
Выполнен ли закон сложения степеней $R^i;R^j=R^{i+j}$, если $i$ и $j$ разного знака?

\end{document}
